\origpage{62}{N. Robertson, P.D. Seymour / Journal of Combinatorial Theory, Series B 89 (2003) 43--76}
\sectionhead{7. The Giant Steps application}

The previous paper of this series contains a result designed to be used at this point. We apply it to prove that 6.1 is implied by the following. (For ``$(\lambda, \mu)$-level'', see [9, section 8].)

\textbf{7.1.} \textit{Let $\Sigma$ be a surface, and let $\tau, \chi, \delta, \lambda \ge 0, \theta' \ge 1$ and $\phi, \psi \ge 3$. Then there exist $\theta \ge 1$ and $\sigma \ge 0$ such that if $\mathcal{T}^*$ is a tangle in a graph $G$, and there is an animal with $\chi$ horns and $\delta$ hairs, of strength $\ge (\theta, \sigma)$, then either}

\hspace*{12pt}$\bullet$ \textit{there is a $\Sigma'$-span of order $\ge \phi$, for some surface $\Sigma'$ obtained from $\Sigma$ by adding a handle, or}

\hspace*{12pt}$\bullet$ \textit{there is an animal with $\chi$ horns and $\delta + 1$ hairs, of strength $\ge (\psi, \tau)$, or}

\hspace*{12pt}$\bullet$ \textit{there exists $Z \subseteq V(G)$ with $|Z| \le \chi + \frac{1}{2}\delta^2\phi^2$ such that some $\Sigma$-span of order $\ge \theta'$ in $G\setminus Z$ with respect to $\mathcal{T}^*\setminus Z$ is $(\lambda, 2\psi)$-level.}

To encourage the reader, let us point out that we are making progress. 7.1 no longer involves segregations of type $(\rho, \kappa)$ nor does it involve controlling minors, which were the two fundamental ingredients of 3.1. Also, crosscaps have gone from the first outcome.

To show that 7.1 implies 6.1 we use the following, which is almost theorem 8.4 of [9].

\textbf{7.2.} \textit{Let $\Sigma$ be a surface and let $p, \phi, \mu \ge 0$. Then there exist $\kappa, \lambda, \rho \ge 0$ and $\theta \ge 1$ such that if $\mathcal{T}^*$ is a tangle in a graph $G$, and some $\Sigma$-span of order $\ge \theta$ is $(\lambda, \mu)$-level, then either:}

\hspace*{12pt}1. \textit{there is a $\Sigma'$-span of order $\phi$, for some surface $\Sigma'$ obtained from $\Sigma$ by adding a crosscap, or}

\hspace*{12pt}2. \textit{there is a $\mathcal{T}^*$-central segregation of $G$ of type $(\rho, \kappa)$ with a proper arrangement in $\Sigma$, or}

\hspace*{12pt}3. \textit{$\mathcal{T}^*$ controls a $K_p$ minor of $G$.}

However, let us point out a discrepancy. We were not farsighted enough in [9], and omitted to include the term ``proper'' in the statement of the theorem, although the \textit{proof} in that paper does yield a proper arrangement. (All the arrangements in [9] first come into being via the proof of theorem 7.7 of that paper, so it is enough to check that they are proper at that stage; and they are, as we can see from theorem 7.6 of [9], or from statement (i) of 7.5 in [9].) Alternately, we could pay for the mistake by starting with a non-proper arrangement, and showing how to convert it to a proper one; that is straightforward, but not very short, and it seems unnecessary to inflict it on the reader.

We also need two other lemmas.
